% status: FULL
\chapter{GPU Architecture for Inference Engineers}\label{app:gpu}
\begin{ChapterIntent}
A GPU is not one very fast arithmetic unit. Its useful abstraction is many
workers sharing several finite storage levels. This chapter prepares the
kernel chapters; generation-specific specifications remain in the dated
frontier ledger and chapters that use them.
\end{ChapterIntent}

\section{Workers, groups and scheduled work}
A kernel launch describes logical workers. In CUDA, threads form blocks
and blocks form a grid; an SM executes resident blocks. Threads in a block
can cooperate through shared memory and synchronization. A CUDA warp
groups 32 threads; other GPU models may use different widths and names
\cite{nvidia2026programmingmodel}. Do not transplant a warp-specific
assumption into portable code without checking it.

The grid describes available work. Occupancy describes resident work
relative to a hardware limit. Neither equals achieved utilization: a
resident worker can wait on memory, a dependency or a barrier. More
occupancy can hide latency but leave less local storage for each useful tile.

\section{Capacity before throughput}
Consider an illustrative SM with 64 KiB of shared memory available to a
kernel and a ceiling of eight resident blocks. If each block needs 24 KiB,
shared memory alone permits
\[
 \left\lfloor 64/24\right\rfloor=2\text{ resident blocks}.
\]
At 16 KiB per block it permits four. Registers, threads and other limits
can reduce either result. This is not a prediction of a twofold speedup:
a smaller tile may require more transfers or less efficient instructions.
Chapter~\ref{ch:matmul-hardware} measures that trade.

\EngineFigure{foundation-memory-levels}{Schematic storage levels, not a
specification or scale drawing. Reusing a local tile avoids device-memory
traffic; keeping weights resident avoids host transfers. Each boundary
needs its own traffic count and bandwidth.}{fig:foundation-memory-levels}

\section{Follow one tile}
Registers hold local values and partial sums. Explicit shared storage holds
reusable tiles for cooperating workers. Caches exploit repeated accesses
without making every reuse explicit in code. Device memory holds large
weight and KV arrays. Spilling registers can add memory traffic; caching
can make repeated source-level loads cheaper than their count suggests.

A matrix tile is loaded, synchronized where required, reused across
products, accumulated and finally stored. The optimization is to replace
expensive repeated movement by cheaper reuse while preserving ordering.
Overwriting a tile before its last reader finishes is a correctness bug,
even if one timing run happens to complete.

Matrix-specialized units accelerate supported matrix operations and
precisions. They do not make exponentials, address translation or tiny
reductions run at the advertised matrix peak. A kernel can hit another
instruction bottleneck while using only a small fraction of matrix FLOP/s.
A useful roofline needs an operation mix as well as a memory boundary.

\section{Capacity, bandwidth and latency}
Capacity says whether an allocation fits. Bandwidth says how many bytes
cross a boundary per second. Latency says how long an individual transfer
or dependency takes. A large device can still have too little bandwidth
for a desired token rate. Many outstanding transfers can hide latency,
but only until another capacity or bandwidth limit binds.

Discrete accelerators add a host-to-device path. An illustrative 64 MiB
transfer at an achieved one-way rate of 16 GiB/s costs at least
\[
 \frac{64\cdot2^{20}}{16\cdot2^{30}}\text{ s}=3.90625\text{ ms}.
\]
Setup and queueing add time. A bidirectional aggregate specification is
not the correct one-way denominator. Overlap hides transfer time only
when dependencies and competing resource demands permit it.

\section{The completion boundary}
A host timer around an asynchronous launch measures submission unless it
waits for the relevant completion. Conversely, synchronizing after every
small kernel can destroy normal overlap. Choose the measurement boundary
before choosing the timer. Chapter~\ref{app:reproducing} supplies a trace.

For AMD and Apple, keep the questions about cooperation, local storage,
residency and completion, but obtain concrete limits from the appropriate
platform documentation. Shared physical memory may remove a host copy;
it does not remove synchronization or make the shared fabric unlimited.
Chapter~\ref{ch:unified-memory} makes that distinction explicit.

\section{Assign work before choosing instructions}
Consider multiplying two small matrices. Each output element is a dot
product of one input row and one input column. A first mapping assigns
one logical worker to each output element. The mapping is easy to explain:
the worker computes its row and column indices, walks the reduction axis,
and writes one result.

This first version can be correct and inefficient. Neighbouring workers
may load many of the same values. A second mapping lets a group cooperate
on a tile of outputs. It loads corresponding input tiles once into
cooperative storage, reuses them for several products, and accumulates
the output tile. The mathematical result has not changed.

There are now two separate shapes to understand. The matrix shape
describes the numerical problem. The launch shape describes how workers
cover that problem. Changing the launch shape must not leave outputs
unwritten or compute boundary elements from invalid memory. Odd
dimensions are useful tests because they expose partially filled tiles.

\section{A barrier protects a dependency}
If one worker loads a value that another worker will use, the consumer
must not read it before the load is available. A group synchronization
can establish that boundary. If the same storage is reused for the
next tile, another dependency prevents overwriting values still being read.

These are correctness requirements, not performance decorations.
Removing a barrier because one small run happened to pass is not an
optimization argument. The execution order may change with input size,
compiler decisions or competing work.

Conversely, unnecessary barriers can force workers to wait when their
work is independent. The useful question is not how many barriers to
remove, but which producer--consumer relationships require ordering.
Draw those relationships for one tile, then choose the synchronization
that the programming model guarantees.

\section{Why more workers can fail to help}
Suppose only a few output tiles are available. Some processors can
have no work even if the kernel is well written. A longer prompt or
a larger batch may create more tiles, allowing more parallelism.
This is a property of the workload shape as well as of the kernel.

At the other extreme, launching many workers does not create more
memory bandwidth. Once useful traffic saturates a shared interface,
additional work can add queueing without increasing the delivered
bytes per second. A memory-bound kernel needs less traffic or better
reuse, not only a larger launch.

Local resource use also matters. A large tile can improve reuse but
consume enough registers or shared memory to reduce the number of
resident groups. The 64 KiB example earlier gives a concrete capacity
calculation. It is only one limit: thread count, registers and other
architecture-specific constraints can bind first.

Run the arithmetic lesson:
\begin{verbatim}
python3 code/foundations/lesson.py gpu
\end{verbatim}
It reports the two storage-limited block counts and the 3.90625 ms
transfer bound. It does not measure a GPU or claim that four resident
blocks run twice as fast as two. Distinguishing a capacity calculation
from a performance result is part of learning the machine.

\section{Matrix units do a particular kind of work}
Matrix-specialized instructions consume tiles in supported layouts and
formats. To benefit, a kernel must supply the right data at the right
time and have enough useful matrix work. A theoretical peak is achieved
under stated assumptions; it is not the rate of every instruction in
the decoder.

Attention also performs reductions, exponentials and address
calculations. Quantized matrix products decode packed values and scales.
A paged cache adds table lookups. These operations can limit a kernel
even when its main matrix instruction is very fast.

Therefore count both the principal arithmetic and the work needed to
feed it. A faster matrix unit can make another part of the same kernel
more visible. This is why later chapters separate an algorithm's FLOP
count from its achieved execution time.

\section{Follow the host and the device separately}
The host prepares arguments and submits work. The device executes that
work later, after any earlier dependencies. The host may return from
submission before the result is ready to read.

A timer stopped at that return answers how long submission took.
A timer stopped after the relevant completion answers a larger
question. Neither measurement is automatically wrong, but labelling
the first as kernel execution time is wrong. State the event that
ends the measurement.

Repeatedly waiting after every small operation can change the workload
by preventing overlap. To time a normal sequence, preserve its
dependency graph and wait at the boundary whose result you actually
need. The launch-overhead chapter measures the cost of different
submission styles on the available hardware.

Graph replay and fusion address different parts of this path.
Replay can reuse a prepared submission structure. Fusion can remove
intermediate launches or data transfers by combining operations.
Both require that the new execution preserve dependencies and
numerical semantics. Neither is a substitute for identifying the
actual bottleneck.

\section{The first GPU experiment should be small}
Before running a model, copy a small known array, execute one operation
and check every result. Increase to an odd-sized array to test the
boundary. Then test the normal synchronization boundary. Only after
those checks should a timing loop run.

Keep the CPU reference even when it is much slower. It can express
the intended computation without the launch geometry or shared-memory
machinery. Compare results before comparing times, and record input
sizes, precision, warm-up and exactly which costs the timer includes.

No GPU is required to complete Part I. The diagrams and arithmetic
explain the execution model, while the existing optional device labs
provide hardware-specific measurements. If the hardware is unavailable,
leave its speed unmeasured. A fabricated peak-to-runtime conversion
would teach the wrong habit.

The important outcome is that you can now point to a worker, a group,
a reusable tile, a memory boundary and a completion event. The
optimized kernel chapters will put real values on those objects.
\section{Worked problems}
\subsection*{1. A residency bound}
\textbf{Problem.} An illustrative processor has 64 KiB available shared
storage and an eight-block ceiling. A block needs 24 KiB. How many fit?
\textbf{Solution.} At most $\min(8,\lfloor64/24\rfloor)=2$ by these limits.
Registers or threads can reduce that count.
\subsection*{2. Transfer time}
\textbf{Problem.} Derive the lower bound for moving 64 MiB at 16 GiB/s.
\textbf{Solution.} The ratio is $64/(16\cdot1024)$ seconds, or 3.90625 ms.
Queueing, setup and other traffic can make the observed time larger.
\subsection*{3. The wrong timer}
\textbf{Problem.} A host call returns before device work completes. What
does a timer around that call establish?
\textbf{Solution.} Submission duration, not completed device execution.
To measure completion, wait for the relevant work or use an appropriate
device timing mechanism with a declared scope.

